\bookchapter{第十九章：故障排查與自救}

\chapterintro{手機同人一樣，偶爾都會「生病」。卡頓、閃退、黑屏、充唔入電……大部分毛病其實冇咁可怕，你完全可以喺送去維修之前自己先試下。呢章就係你嘅「手機急診手冊」——先用低成本嘅方法試，唔得再揾專業醫生。}

\booksection{一、手機變卡、發熱——最常被問嘅問題}

講真，手機用咗一年左右開始變卡，太正常喇。原因唔一定係手機壞咗，更多時候係「積灰」——垃圾檔案、後台進程、存儲空間慢慢被填滿。

\booksubsection{先做診斷，唔好急住修：}

\begin{enumerate}
  \item \textbf{查存儲空間}：去設定 存儲睇剩低空間。如果唔夠10\%喇，系統會變得非常卡，因為佢喺掙扎住騰地方。我嘅經驗係：保持至少15\%-20\%嘅剩低空間，手機先可以比較順暢。
  \item \textbf{查耗電排行}：去電池設定入面睇過去24個鐘嘅耗電排行。如果某個App佔咗大頭但你根本冇點用，咁就係佢喺後台偷跑。直接卸載或者限制後台活動。
  \item \textbf{聯想最近做咗乜}：係升級咗系統後變卡嘅？（大版本更新後前幾日變卡正常，系統喺後台重建索引）。仲係安裝咗某個新App之後變卡嘅？（卸載試下）。
  \item \textbf{睇使用環境}：手機喺太陽底下暴曬、套住厚厚嘅殼邊充邊打機，發燙係必然嘅。高温仲會導致CPU自動降頻，你就感覺手機突然變卡喇。涼返再試。
\end{enumerate}

\textbf{治療方法（由簡單到複雜）}：重啟手機（解決60\%嘅問題） 清存儲空間（刪掉冇用嘅App同影片） 某個App有問題就清佢嘅快取或卸載重裝 最後一招：備份數據後恢復出廠設定。

\booksection{二、應用閃退、無響應——App唔聽話喇}

某個App撳開就閃退、或者用着用着突然卡死唔鬱。處理流程就好似醫生問診，一步步嚟：

\begin{enumerate}
  \item \textbf{先關掉重開}：由後台劃走App，再重新打開。大多數臨時bug呢一下就好喇。
  \item \textbf{更新App}：去應用商店睇有冇更新。有時候係因為App版本太舊，同系統唔兼容。
  \item \textbf{檢查權限}：有啲App需要某個權限先可以正常工作，但你冇畀。比如相機App冇畀相機權限打開就閃退。去設定入面揾到嗰個App，將該畀嘅權限都畀咗。
  \item \textbf{清快取}：去設定 應用管理 揾到嗰個App 存儲 清除快取。注意係「清除快取」唔係「清除數據」！清除數據會將你App入面嘅登入狀態同設定全刪咗，相當於恢復出廠畀App嘅狀態。
  \item \textbf{卸載重裝}：最後一招。卸載之前先確認呢個App嘅數據會唔會冇——微信嘅話先備份聊天記錄，遊戲嘅話先確認帳號有冇綁定手機號。
\end{enumerate}

\booksection{三、充唔入電或充電慢——可能唔係電池嘅鍋}

充唔入電係特別容易等人焦慮嘅問題，但好多時候原因係出喺配件上，唔係手機壞咗。

\begin{itemize}
  \item \textbf{先排查充電口}：用牙籤或者唔導電嘅細工具，輕輕咁將充電口入面嘅灰塵同棉絮掏一掏。呢個我自己遇過——口袋入面跌咗一條線頭啱啱堵喺充電口入面，插落去以為係接觸不良，掏出嚟就好喇。
  \item \textbf{換根線同充電頭試下}：線嘅損壞率比充電頭高得多，尤其係經常卷來捲去嘅線。換一條確定冇問題嘅線同充電頭交叉測試，好快可以確定係邊個配件壞咗。
  \item \textbf{無線充電嘅問題}：如果你用無線充電，檢查一下手機殼係咪太厚，或者殼後面有冇金屬片（磁吸支架之類）。仲有，將手機準確咁放到充電板嘅正中間——歪咗就可能充唔上。
  \item \textbf{充電慢}：確認充電頭功率夠唔夠（iPhone想要快充至少需要20W嘅PD充電頭），線支唔支援快充（唔係所有嘅線都可以過3A以上電流），手機係咪太燙（太燙會自動降功率保護電池）。
\end{itemize}

\booksection{四、網絡連唔上、Wi-Fi成日斷}

喺出面用流量好地地，返屋企連唔上Wi-Fi；或者Wi-Fi連住但就係冇網。按呢個順序排查：

\begin{enumerate}
  \item \textbf{開關飛行模式}：下拉控制中心，打開飛行模式等5秒再關掉。呢個等於等手機重新搜一次網，非常有效。
  \item \textbf{Wi-Fi「忘記」再重連}：去Wi-Fi設定入面，撳嗰個網絡名旁邊嘅感歎號（安卓係齒輪），揀「忽略此網絡」或「忘記」，然後重新輸入密碼連接。
  \item \textbf{重啟路由器}：好多人屋企嘅路由器幾個月冇關過，時間耐咗會唔太穩定。拔掉電源等30秒再插返去，大部分Wi-Fi問題就解決喇。
  \item \textbf{檢查係咪手機嘅問題}：別嘅設備連同一個Wi-Fi正常嗎？如果都正常就你嘅手機唔得，咁手機可能需要重啟或者重置網絡設定（設定 通用 傳輸或還原 還原網絡設定，唔會刪你嘅數據，但會將所有Wi-Fi密碼都清咗）。
\end{enumerate}

\booksection{五、螢幕問題：觸控唔靈、閃屏、留印子}

\begin{itemize}
  \item \textbf{觸控唔靈敏}：先撕掉鋼化膜睇下係咪膜嘅鍋——平鋼化膜用耐咗確實會影響觸控。螢幕上有水或手上有油都會導致觸控漂移。都唔係嘅話，重啟一下再試。
  \item \textbf{螢幕上有印子/殘影（燒屏）}：長時間顯示同一個畫面（比如一直亮屏碌抖音），會喺OLED螢幕上留低「殘影」。呢個唔係螢幕壞咗，淨係暫時嘅。關屏休息半個鐘一般可以恢復。
  \item \textbf{螢幕閃爍}：如果係低亮度下閃爍，可能係PWM調光嘅問題，多見於OLED螢幕。去設定入面揾「DC調光」或「防頻閃」選項，打開後會好好多。唔係所有手機都有呢個功能。
  \item \textbf{螢幕有一塊區域完全黑咗/亮咗}：呢個比較嚴重，大概率係螢幕物理損壞，需要換屏。非官方維修換嘅螢幕，色彩同亮度通常不如原裝，都影響防水。
\end{itemize}

\booksection{六、無法開機、卡喺Logo介面——先唔好慌}

呢種情況最容易等人心跳加速，但冷靜，先試下呢幾個：

\begin{itemize}
  \item \textbf{充電半個鐘再講}：好多時候淨係電徹底耗盡咗。插上原裝充電器，等至少30分鐘，再嘗試開機。有啲手機喺電耗盡後需要充到一定電量先會有反應。
  \item \textbf{強制重啟}：iPhone全面屏機型撳一下音量加、一下音量減、長撳側邊電源鍵唔放，直到螢幕出現蘋果Logo（可能要撳住10-15秒）。安卓通用方法：長撳電源鍵10秒以上，直到震動或亮屏。
  \item \textbf{卡喺啟動介面反覆重啟}：呢種情況叫「白蘋果」（iPhone）或「卡Logo」（安卓）。如果強制重啟唔得，可能需要進入恢復模式用電腦修復。iPhone連電腦打開iTunes/Finder，進入恢復模式後揀「更新」（唔係「恢復」），會重裝系統但保留數據。安卓進Fastboot或Recovery模式，各品牌組合鍵唔同，建議搜尋你手機型號+「進恢復模式」。
  \item \textbf{以上都唔得}：那就只能求助專業維修喇。數據如果備份過，就唔使太擔心。
\end{itemize}

\booksection{七、聲音問題：喇叭唔響、通話聲音細}

\begin{itemize}
  \item \textbf{喇叭冇聲}：先確定唔係靜音模式（iPhone側邊嘅靜音撥片）、勿擾模式、或者音量被撳到最低咗。聽落好似好基本，但呢種情況佔咗40\%。然後檢查藍牙係咪連住某個耳機或音箱。
  \item \textbf{通話對方聽唔到你講嘢}：檢查話筒（麥克風）孔有冇堵住——手機底部嘅細孔好容易俾灰塵堵住。用軟毛刷輕輕清理。唔好用針或牙籤硬拮，會拮壞防塵網。
  \item \textbf{聲音沙啞/破音}：如果係喇叭入水導致嘅——唔好用風筒熱風吹！用紙巾吸乾表面水分，然後放住自然晾乾。網上話用聲波排水App有一定效果（播放低頻聲音將水振出去），可以試下。
\end{itemize}

\booksection{八、揾官方維修，避開嗰啲坑}

如果上面嘅方法都試過都唔得，咁就要揾專業維修喇。但維修市場水好深，講下我自己踩過嘅坑同經驗：

\begin{itemize}
  \item \textbf{揾真官方，唔好揾「睇落似官方」嘅}：去品牌官網（Apple.com.cn、Huawei.com.cn等）揾「授權服務商」列表。街邊掛個蘋果Logo嘅非授權店，用嘅可能唔係原裝配件，修完你嘅防水就冇咗一剛。
  \item \textbf{送修前必須做嘅事}：備份所有數據（唔好指望修嘅時候數據唔會冇），退出所有重要帳號（Apple ID、微信、支付寶），關閉「查找我嘅手機」。有啲維修點會要求你提供鎖屏密碼，呢件事有私隱風險——你最好自己先退出帳號再送過去。
  \item \textbf{換螢幕嘅坑}：非官方螢幕分「原壓屏」（原裝玻璃換過）、「後壓屏」（非原裝玻璃）、「組裝屏」（全非原裝）。價格差好幾倍，但顯示效果同耐久度都差好多。如果手機仲喺保修期，老老實實走官方。
  \item \textbf{換電池嘅坑}：非原裝電池可能虛標容量，而且用咗半年就衰減得比原裝電池兩年嘅損耗仲大。蘋果換電池官方價格比非官方貴唔了多少，建議走官方。
  \item \textbf{數據恢復嘅坑}：手機入水/摔壞後，去維修如果對方話「數據恢復要加錢」，小心——有啲不良商家會故意誇大損壞程度嚟多收錢。正規嘅數據恢復係先報價再操作，而且恢復唔成功唔收費。
\end{itemize}

\begin{tcbwarn}{萬能第一招：重啟手機}
唔管你遇到乜嘢問題——卡頓、閃退、連唔上網、聲音異常——先重啟一下手機。呢句話聽落似係敷衍，但以我自己嘅經驗，至少有六成嘅小毛病重啟之後就好喇。手機同電腦一樣，跑得耐記憶體入面會積累各種亂七八糟嘅嘢，重啟就係清空重來。所以，下次手機出問題，先唔好急住搜「點算」，重啟完咗再睇下。

\end{tcbwarn}

\begin{tcbpractice}{而家就記：呢啲救命操作要爛熟於心}
\begin{enumerate}
  \item \textbf{強制重啟}：全面屏iPhone撳音量加 音量減 長撳電源鍵到出現Logo。安卓長撳電源鍵10-15秒直到震動。
  \item \textbf{清理充電口}：揾一個細工具（牙籤或手機取卡針），輕輕清理充電口嘅灰塵。
  \item \textbf{Wi-Fi重連}：去Wi-Fi設定入面「忘記」你當前嘅Wi-Fi，重新輸入密碼連接。
  \item \textbf{開/關飛行模式}：下拉控制中心，打開飛行模式，心入面數5秒，關掉，睇網絡係咪恢復。
  \item 將呢幾個操作嘅步驟截圖保存，或者寫喺備忘錄入面。遇到問題嘅時候唔使慌慌張張去搜，直接揭備忘錄就得。
\end{enumerate}

\end{tcbpractice}
